% gw_predictions_section.tex
% Quantitative Gravitational Wave Predictions from Sigma = 2 ln Q
% Supplement for Paper 2
% Author: Sheng-Kai Huang
% Compile: pdflatex gw_predictions_section.tex

\documentclass[prl,twocolumn,superscriptaddress,longbibliography]{revtex4-2}

\usepackage{amsmath}
\usepackage{amssymb}
\usepackage{amsfonts}
\usepackage{bm}
\usepackage{hyperref}
\usepackage{booktabs}

% Macros
\newcommand{\kB}{k_{\mathrm{B}}}
\newcommand{\Tr}{\mathrm{Tr}}
\newcommand{\Petz}{\mathcal{R}}
\newcommand{\chan}{\mathcal{N}}
\newcommand{\ket}[1]{|#1\rangle}
\newcommand{\bra}[1]{\langle#1|}
\newcommand{\tauasym}{\tau}
\newcommand{\ceff}{c_{\mathrm{eff}}}
\newcommand{\rs}{r_{\mathrm{s}}}
\newcommand{\Sgrav}{\Sigma_{\mathrm{grav}}}
\newcommand{\Fbound}{F_{\mathrm{bound}}}
\newcommand{\nref}{n_{\mathrm{grav}}}
\newcommand{\cs}{c_{\mathrm{s}}}

\begin{document}

\title{Gravitational Wave Predictions of the $\Sigma = 2\ln Q$ Framework:\\
Quantitative Estimates for Current and Near-Future Detectors}

\author{Sheng-Kai Huang}
\email{akai@fawstudio.com}
\affiliation{Independent Researcher}

\date{\today}

\begin{abstract}
We derive three specific, testable gravitational wave predictions
from the $\Sigma = 2\ln Q$ framework developed in the companion
papers.  (1)~The Khronon ghost condensation mechanism predicts
\emph{no detectable scalar gravitational wave mode} in any current
or planned detector: the Jeans frequency $f_J \sim 10^{-19}$~Hz
lies far below all detector bands, and the DBI-saturated scalar
strain is bounded by $h_s < \alpha_{\mathrm{DBI}}\,h_t < 0.05\,h_t$.
(2)~The exponential metric $g_{00} = -e^{-\rs/r}$, derived from
Petz bound saturation, produces a cumulative phase shift of
$\Delta\Phi \sim 10^4$--$10^6$~radians in extreme-mass-ratio
inspirals (EMRIs) observable by LISA, with a minimum detectable
fractional deviation from Schwarzschild of $\alpha_{\min}
\sim 10^{-6}$.  (3)~The framework predicts holographic scaling
($\alpha_{\mathrm{dec}} = 0$) for gravitational wave phase
diffusion, distinguishable from string foam ($\alpha_{\mathrm{dec}}
= -1$) and causal set ($\alpha_{\mathrm{dec}} = +1$) predictions
by the Einstein Telescope and Cosmic Explorer.
\end{abstract}

\maketitle

%=======================================================================
\section{Introduction}
%=======================================================================

The $\Sigma = 2\ln Q$ framework~\cite{Paper1,Paper2,Paper3}
identifies gravitational entropy production with the quantum
relative entropy (QRE) and predicts the exponential metric as
the Petz-saturated geometry.  While Paper~2 established three
strong-field predictions (shadow size, QNM frequency shift,
and echoes), several gravitational wave signatures remain to
be quantified.  Here we provide explicit numerical predictions
that can be confronted with data from current LIGO--Virgo--KAGRA
observations, the upcoming LISA mission, and next-generation
detectors (Einstein Telescope, Cosmic Explorer).

%=======================================================================
\section{Prediction 1: No Scalar Gravitational Wave Mode}
\label{sec:scalar}
%=======================================================================

\subsection{Khronon perturbation sound speed}

In the Blanchet--Skordis Khronon theory~\cite{BS2024,BS2025},
the scalar field $\phi$ sits in a ghost condensation background
with kinetic function $K(Q)$, where $Q = |\partial\phi|$ is the
inverse Khronon lapse.  The perturbation sound speed is
\begin{equation}
  \cs^2 = \frac{K'(Q_0)}{K'(Q_0) + 2\,Q_0\,K''(Q_0)}
  \;\xrightarrow{K'(Q_0) = 0}\; 0\,.
  \label{eq:cs2}
\end{equation}
Ghost condensation ($K'(Q_0) = 0$) gives $\cs^2 = 0$
exactly~\cite{AHCLM2004}, which is radiatively stable with
quantum corrections $\delta\cs^2 \sim (\mu/M_{\mathrm{Pl}})^2
\sim 10^{-122}$~\cite{Paper3}.

The DBI completion~\cite{BS2024} introduces a scale-dependent
correction:
\begin{equation}
  \cs^2(k,a) = \alpha_{\mathrm{DBI}}\,
  \frac{(k/k_J)^2}{1 + (k/k_J)^2}\,,
  \label{eq:cs2_DBI}
\end{equation}
where $\alpha_{\mathrm{DBI}} = Q_0/(2\lambda_D^2)$ and the
Jeans wavenumber is
\begin{equation}
  k_J(a) = \sqrt{\tfrac{3}{2}\,\Omega_K\,H_0^2/a}\,,
  \label{eq:kJ}
\end{equation}
with $\Omega_K \approx 0.265$ the Khronon energy fraction.

\subsection{Jeans frequency}

Evaluating at $a = 1$ (today) with $H_0 = 67.36$~km/s/Mpc:
\begin{equation}
  k_J = 1.42 \times 10^{-4}\;\mathrm{Mpc}^{-1}
  = 4.59 \times 10^{-27}\;\mathrm{m}^{-1}\,.
\end{equation}
The corresponding frequency is
\begin{equation}
  \boxed{f_J = \frac{c\,k_J}{2\pi}
  = 2.2 \times 10^{-19}\;\mathrm{Hz}}\,,
  \label{eq:fJ}
\end{equation}
which is $10^{10}$ times below nanohertz pulsar timing
arrays and $10^{17}$ times below the LIGO band.

At all detector frequencies ($f \gg f_J$), the ratio $k/k_J
\gg 1$ and $\cs^2 \to \alpha_{\mathrm{DBI}}$.  The CMB
constraint from Paper~3 gives $\alpha_{\mathrm{DBI}} < 0.05$;
the ghost condensation natural value is $\alpha_{\mathrm{DBI}}
\lesssim 0.01$.

\subsection{Scalar mode amplitude}

The scalar gravitational wave strain couples to the detector
through the breathing mode polarization.  In the
Blanchet--Skordis theory, the scalar-to-tensor amplitude ratio
is bounded by~\cite{BS2024,SchumacherEA2024}
\begin{equation}
  \frac{h_{\mathrm{scalar}}}{h_{\mathrm{tensor}}}
  \lesssim \alpha_{\mathrm{DBI}} < 0.05\,.
\end{equation}
At Advanced LIGO sensitivity $h \sim 10^{-23}$, the scalar mode
is $h_s \lesssim 5 \times 10^{-25}$, well below the noise floor.

\emph{Prediction.}---No scalar gravitational wave mode is
detectable by any current or planned detector (LIGO, LISA, ET,
CE, PTA).  The detection of a propagating scalar mode at any
frequency would falsify the ghost condensation mechanism
$K'(Q_0) = 0$ that underpins the framework.

%=======================================================================
\section{Prediction 2: EMRI Phase Shift from Exponential Metric}
\label{sec:emri}
%=======================================================================

\subsection{Metric comparison at the same areal radius}

The exponential metric in isotropic coordinates~\cite{Paper2},
\begin{equation}
  ds^2 = -e^{-\rs/r_{\mathrm{iso}}}\,c^2\,dt^2
       + e^{+\rs/r_{\mathrm{iso}}}
       \!\left(dr_{\mathrm{iso}}^2
       + r_{\mathrm{iso}}^2\,d\Omega^2\right),
  \label{eq:exp_metric}
\end{equation}
must be compared with Schwarzschild at the same
circumferential (areal) radius $R$.  The coordinate
transformation $R = r_{\mathrm{iso}}\,
\exp\!\big(GM/(c^2 r_{\mathrm{iso}})\big)$ is inverted
numerically.

At the same areal radius $R = u\,GM/c^2$, the metric
components differ by
\begin{equation}
  \frac{\Delta g_{00}}{g_{00}}
  = \frac{g_{00}^{\mathrm{exp}}(R) - g_{00}^{\mathrm{Schw}}(R)}
         {g_{00}^{\mathrm{Schw}}(R)}
  \sim O\!\left(\frac{GM}{c^2 R}\right)^{\!3}\,,
  \label{eq:delta_g00}
\end{equation}
consistent with the $O((\rs/r)^3)$ first discrepancy noted
in Paper~2.  Table~\ref{tab:metric} gives numerical values.

\begin{table}[t]
\caption{\label{tab:metric}%
Comparison of the exponential and Schwarzschild metrics at the
same areal radius $R = u\,GM/c^2$.  The fractional orbital
frequency difference $\delta\omega/\omega$ is computed from the
effective potential.}
\begin{ruledtabular}
\begin{tabular}{rccc}
  $u$ & $g_{00}^{\mathrm{Schw}}$ & $g_{00}^{\mathrm{exp}}$
    & $\delta\omega/\omega$ \\
  \colrule
  6   & $-0.6667$ & $-0.6643$ & $+1.22\%$ \\
  8   & $-0.7500$ & $-0.7491$ & $+0.58\%$ \\
  10  & $-0.8000$ & $-0.7996$ & $+0.34\%$ \\
  15  & $-0.8667$ & $-0.8666$ & $+0.13\%$ \\
  20  & $-0.9000$ & $-0.9000$ & $+0.07\%$ \\
  50  & $-0.9600$ & $-0.9600$ & $+0.001\%$ \\
\end{tabular}
\end{ruledtabular}
\end{table}

\subsection{ISCO comparison}

The innermost stable circular orbit (ISCO) is determined by
the vanishing of $d^2 V_{\mathrm{eff}}/dR^2$.  Numerical
evaluation gives:
\begin{equation}
  R_{\mathrm{ISCO}}^{\mathrm{exp}} = 6.34\,\frac{GM}{c^2}\,,
  \qquad
  R_{\mathrm{ISCO}}^{\mathrm{Schw}} = 6.00\,\frac{GM}{c^2}\,,
\end{equation}
a $5.7\%$ outward shift.  This modifies the final plunge
waveform and the gravitational wave cutoff frequency.

\subsection{Accumulated EMRI phase difference}

For an extreme-mass-ratio inspiral with mass ratio
$q = m/M \ll 1$, the number of orbits from radius $R$ to
$R_{\mathrm{ISCO}}$ in the adiabatic approximation is
\begin{equation}
  N_{\mathrm{orb}}(u)
  = \frac{1}{64\pi q}
  \left(u^{5/2} - u_{\mathrm{ISCO}}^{5/2}\right),
\end{equation}
where $u = R c^2/(GM)$.  The accumulated phase difference
between the exponential and Schwarzschild waveforms is
\begin{equation}
  \Delta\Phi = \frac{5}{32\,q}\int_{u_{\mathrm{ISCO}}}^{u_{\max}}
  \!\frac{\delta\omega}{\omega}(u)\;u^{3/2}\,du\,,
  \label{eq:DeltaPhi}
\end{equation}
where $\delta\omega/\omega$ is the fractional orbital frequency
difference at areal radius $u\,GM/c^2$.

Numerical integration from $u_{\mathrm{ISCO}} = 6$ to
$u_{\max} = 100$ gives, for representative EMRI parameters:

\begin{table}[t]
\caption{\label{tab:emri}%
Accumulated GW phase difference between exponential and
Schwarzschild metrics for EMRIs, computed via
Eq.~\eqref{eq:DeltaPhi}.  LISA phase sensitivity is
$\sigma_\Phi \sim 0.1$~rad.}
\begin{ruledtabular}
\begin{tabular}{lcccc}
  $M/M_\odot$ & $m/M_\odot$ & $q$ & $\Delta\Phi$~(rad)
    & $f_{\mathrm{ISCO}}$~(mHz) \\
  \colrule
  $10^5$ & 10 & $10^{-4}$ & $6.7\times 10^3$ & 43.8 \\
  $10^6$ & 10 & $10^{-5}$ & $6.7\times 10^4$ & 4.38 \\
  $10^6$ & 30 & $3\times 10^{-5}$ & $2.2\times 10^4$ & 4.38 \\
  $10^7$ & 10 & $10^{-6}$ & $6.7\times 10^5$ & 0.44 \\
\end{tabular}
\end{ruledtabular}
\end{table}

\emph{Prediction.}---For a standard EMRI ($M = 10^6\,M_\odot$,
$m = 10\,M_\odot$):
\begin{equation}
  \boxed{\Delta\Phi \approx 6.7 \times 10^4\;\mathrm{rad}}\,.
  \label{eq:DeltaPhi_value}
\end{equation}
With LISA phase sensitivity $\sigma_\Phi \sim 0.1$~rad,
this corresponds to a detection significance of
$\sim 7 \times 10^5\,\sigma$.  Even if the exponential metric
contributes only as a fractional correction $\alpha$ to the
Schwarzschild geometry, LISA can detect deviations as small as
\begin{equation}
  \alpha_{\min} = \frac{\sigma_\Phi}{\Delta\Phi}
  \approx 1.5 \times 10^{-6}\,.
  \label{eq:alpha_min}
\end{equation}

\emph{Comparison with LQG predictions.}---Zi \& Kumar~\cite{ZiKumar2024}
found that LISA can detect loop quantum gravity modifications
in EMRIs at the $2 \times 10^{-6}$ fractional level.  The
exponential metric prediction~\eqref{eq:DeltaPhi_value}
is $\sim 10^5$ times larger than the LQG effect, making it
far easier to detect---or exclude.

\subsection{Consistency with current LIGO observations}

LIGO observes the final $\sim$10--30 orbits of comparable-mass
binary mergers at $r \sim 3$--$5\,\rs$.  At these radii,
$\delta\omega/\omega \sim 1$--$2\%$.  Over $\sim 20$ orbits,
the accumulated phase difference is $\Delta\Phi \sim 1$--$4$~rad.
Current systematic uncertainties in LIGO waveform models are
$\sim 1$~rad~\cite{GWTC4_TGR2026}, so the exponential metric
is \emph{not yet excluded} by LIGO observations but lies at
the boundary of detectability.  The GWTC-4 analysis
finding consistency with GR constrains the exponential metric
admixture to $\alpha \lesssim 0.1$--$1$ (order-of-magnitude
estimate).  Next-generation detectors (ET, CE) with
$\sim 10\times$ better SNR will provide definitive tests.

%=======================================================================
\section{Prediction 3: Gravitational Wave Phase Diffusion}
\label{sec:decoherence}
%=======================================================================

\subsection{GW decoherence from the Sigma framework}

From Paper~2, the Pikovski gravitational decoherence rate
in the $\tauasym$ framework is
\begin{equation}
  \tauasym_{\mathrm{grav}}(t)
  = 1 - \exp\!\left(-N\,\Gamma\,t^2\right),
\end{equation}
with $\Gamma = (\Delta E)^2\,(\Delta\Phi)^2 / (2\hbar^2 c^4)$.
In terms of $\Sigma$:
$\Sgrav = 2N\Gamma t^2$.

A gravitational wave with strain $h$ and angular frequency
$\omega_{\mathrm{GW}}$ produces a tidal acceleration
$g_{\mathrm{eff}} = h\,\omega_{\mathrm{GW}}^2\,\Delta x$
across a quantum superposition of size $\Delta x$, giving
a decoherence timescale
\begin{equation}
  t_{\mathrm{dec}} = \frac{\hbar}{m\,g_{\mathrm{eff}}\,\Delta x}
  = \frac{\hbar}{m\,h\,\omega_{\mathrm{GW}}^2\,\Delta x^2}\,.
  \label{eq:t_dec}
\end{equation}

For a nanoparticle interference experiment
($m = 170$~kDa~\cite{Pedalino2026},
$\Delta x = 1\;\mu$m) and a GW event
($h = 10^{-21}$, $f = 100$~Hz, $\Delta t = 0.1$~s):
\begin{align}
  t_{\mathrm{dec}} &\approx 9.5 \times 10^{14}\;\mathrm{s}\,,
  \nonumber \\
  \Sgrav &= 2\,\Delta t / t_{\mathrm{dec}}
  \approx 2.1 \times 10^{-16}\,.
  \label{eq:Sigma_GW}
\end{align}
Individual gravitational wave events produce completely
negligible decoherence in any realistic quantum system.

\subsection{Phase diffusion exponent}

The $\Sigma$ framework makes a more consequential prediction for
the \emph{propagation} decoherence of gravitational waves
themselves.  Cang \& Wang~\cite{CangWang2025} showed that
quantum spacetime fluctuations produce phase diffusion with
variance
\begin{equation}
  \mathrm{var}(\phi) \propto d \cdot f^{\alpha_{\mathrm{dec}}}\,,
  \label{eq:phase_diffusion}
\end{equation}
where $d$ is the propagation distance and the frequency exponent
$\alpha_{\mathrm{dec}}$ discriminates between quantum gravity
models.

In the $\Sigma$ framework, information loss accumulates as
$\Sgrav = D(\rho \| \sigma)$ along the propagation path.
Since the QRE is additive (extensivity, Paper~1) and
frequency-independent at leading order, the phase variance
scales as $d$ with no frequency dependence:
\begin{equation}
  \boxed{\alpha_{\mathrm{dec}} = 0
  \quad(\text{holographic scaling})}\,.
  \label{eq:alpha_dec}
\end{equation}

This is distinguished from competing models:
\begin{center}
\begin{tabular}{lc}
\toprule
Model & $\alpha_{\mathrm{dec}}$ \\
\midrule
String foam & $-1$ \\
$\Sigma = 2\ln Q$ (this work) & $0$ \\
Causal set & $+1$ \\
\bottomrule
\end{tabular}
\end{center}

\emph{Prediction.}---Current bounds from
Planck + BICEP/Keck + LVK + BBN~\cite{deKruijf2025} are
consistent with $\alpha_{\mathrm{dec}} = 0$.  The Einstein
Telescope and Cosmic Explorer will improve sensitivity by
$\sim 100\times$, enabling a definitive measurement of
$\alpha_{\mathrm{dec}}$.

%=======================================================================
\section{Summary of Falsifiable Predictions}
\label{sec:summary}
%=======================================================================

Table~\ref{tab:summary} collects all quantitative predictions
with specific numerical values and detector requirements.

\begin{table*}[t]
\caption{\label{tab:summary}%
Falsifiable gravitational wave predictions of the
$\Sigma = 2\ln Q$ framework.  All predictions have specific
numerical values testable by identified experiments.}
\begin{ruledtabular}
\begin{tabular}{llccl}
  \textbf{Prediction} & \textbf{Value} & \textbf{Detector}
    & \textbf{Sensitivity} & \textbf{Timeline} \\
  \colrule
  No scalar GW mode & $h_s < 0.05\,h_t$ & LIGO/ET/CE
    & $h \sim 10^{-23}$ & Now (null result) \\
  EMRI phase shift & $\sim 10^4$--$10^6$~rad & LISA
    & 0.1~rad & 2035+ \\
  GW echoes (stellar BH) & $\Delta t = 2.5$~ms
    & LIGO O5 & Need templates & 2025--2027 \\
  GW echoes (SMBH) & $\Delta t = 165$~s
    & LISA & Standard search & 2035+ \\
  QNM frequency shift & $-4.4\%$
    & ET/CE & $\sim 1\%$ & 2035+ \\
  Shadow size & $+4.6\%$ & ngEHT & $\sim 2$--$3\%$ & 2030+ \\
  Phase diffusion exponent & $\alpha_{\mathrm{dec}} = 0$
    & ET/CE & $100\times$ improvement & 2035+ \\
  GW friction $\Xi_0$ & $= 1$ exactly
    & Standard sirens & $10^{-3}$ & 2035+ \\
  ISCO shift & $+5.7\%$ & LISA (EMRI) & $\sim 1\%$ & 2035+ \\
\end{tabular}
\end{ruledtabular}
\end{table*}

The most discriminating near-term test is the search for
gravitational wave echoes at $\sim 2.5$~ms delay in existing
LIGO--Virgo--KAGRA data using short-delay templates.  The most
powerful future test is the EMRI phase shift measurable by LISA,
which can detect deviations from Schwarzschild at the
$\alpha \sim 10^{-6}$ level.

%=======================================================================
% REFERENCES
%=======================================================================

\begin{thebibliography}{99}

\bibitem{Paper1}
S.-K.~Huang, ``Petz recovery unification'' (Paper~I).

\bibitem{Paper2}
S.-K.~Huang, ``The gravitational refractive index as Petz
recovery fidelity'' (Paper~II).

\bibitem{Paper3}
S.-K.~Huang, ``Weak-field Khronon dark matter from
$\Sigma = 2\ln Q$'' (Paper~III).

\bibitem{BS2024}
L.~Blanchet and C.~Skordis,
``Relativistic Khronon Theory in agreement with Modified
Newtonian Dynamics and Large-Scale Cosmology,''
JCAP \textbf{11}, 040 (2024);
arXiv:2404.06584.

\bibitem{BS2025}
L.~Blanchet and C.~Skordis,
``Khronon-Tensor theory reproducing MOND and the cosmological
model,''
arXiv:2507.00912 (2025).

\bibitem{AHCLM2004}
N.~Arkani-Hamed, H.-C.~Cheng, M.~A.~Luty, S.~Mukohyama,
and T.~Wiseman,
``Ghost condensation and a consistent infrared modification
of gravity,''
JHEP \textbf{05}, 074 (2004).

\bibitem{SchumacherEA2024}
A.~M.~Schumacher, N.~Yunes, \emph{et~al.},
``Gravitational wave constraints on Einstein-aether theory
with LIGO/Virgo data,''
Phys.\ Rev.\ D \textbf{108}, 104046 (2024);
arXiv:2304.06801.

\bibitem{ZiKumar2024}
T.~Zi and S.~Kumar,
``Eccentric extreme mass-ratio inspirals: A gateway to probe
quantum gravity effects,''
arXiv:2409.17765 (2024).

\bibitem{GWTC4_TGR2026}
LIGO--Virgo--KAGRA Collaboration,
``GWTC-4.0: Tests of General Relativity. I,''
arXiv:2603.19019 (2026).

\bibitem{Pedalino2026}
F.~Pedalino \emph{et~al.},
``Quantum interference of a 170 kDa nanoparticle,''
Nature \textbf{649}, 866 (2026).

\bibitem{CangWang2025}
J.~Cang and Y.-T.~Wang,
``Universality and Falsifiability of Quantum Spacetime
Decoherence,''
arXiv:2512.02782 (2025).

\bibitem{deKruijf2025}
J.~de Kruijf, C.~Galloni, and N.~Bartolo,
``The first data-driven bounds on the quantum decoherence
of inflationary gravitational waves,''
arXiv:2511.14727 (2025).

\bibitem{Cardoso2016}
V.~Cardoso, E.~Franzosi, and P.~Pani,
``Is the gravitational-wave ringdown a probe of the event
horizon?''
Phys.\ Rev.\ Lett.\ \textbf{116}, 171101 (2016).

\end{thebibliography}

\end{document}
